% Appendix companion to Section 3. Procedure and validity evidence.
% Dataset: dataset/multi_domain_dataset/.

\subsection{Dataset Construction Pipeline}
\label{app:dataset}

\paragraph{Sources and idea extraction.}
Papers are drawn from arXiv in three areas: graph machine learning, computer
vision, and time series. Selection prefers work that was subsequently published
and, among that, papers at major venues. From each paper's title, abstract and
introduction an LLM extracts
innovation claims and assigns each an innovation aspect; a second LLM pass
discards claims that are not innovation claims at all, including open-source
and reproducibility statements, bare state-of-the-art assertions, and
engineering deliverables.

\paragraph{Candidate priors.}
Idea texts are embedded and indexed for top-$K$ nearest-neighbour recall, and
candidate relations are formed under a hybrid rule: high similarity is accepted
directly, low similarity is rejected directly, and the uncertain band is passed
to an LLM. The base build pairs ideas of the same type only, which would record
an idea whose closest relative is of another type as isolated. A cross-type
pass recovers these relations.
The retained priors are ranked for each target. Pairing stays within a source
domain, and no cross-domain pair is labelled, so the three single-domain subsets
and their union address the same ideas under identifiers namespaced by domain.
The assembled dataset passes its own validation checks with no duplicate pair
keys, no duplicate novelty targets, no invalid pair labels, no target left
without a pair, and no temporal violation.

\paragraph{Annotation.}
All annotation stages use the same judge, \texttt{deepseek-v4-pro}. It labels
pair coverage over the four classes. Pairs in the ambiguous band are
re-judged rather than accepted from the similarity rule. This band includes
related but not covering, partial cover, large cover, and high-similarity
pairs the graph left unconnected. It then
estimates, for each target, the substantive residual that survives the retained
evidence, the sufficiency of that evidence, and the extent to which the claimed
difference is surface-level. Idea-level grades follow by applying
the novelty rule of Sec.~\ref{sec:dataset-task} to those quantities. Labels are model-generated;
each stage was reviewed by sampled manual inspection.

\paragraph{Composition.}
\textsc{IdeaDelta} contains $21{,}111$ annotated ideas and $92{,}876$ annotated
target--prior pairs, with target dates from 2019-01 to 2026-09
(Table~\ref{tab:dataset-composition}). Each idea carries its source domain, and
the \emph{All} column is the union of the three domain subsets rather than a
fourth collection.

\begin{table}[t]
\centering
\small
\setlength{\tabcolsep}{6pt}
\caption{Composition of \textsc{IdeaDelta}: counts of ideas and of
target--prior pairs, and the idea-level novelty distribution, by source domain.
Pairs are formed within a source domain only, so \emph{All} is the union of the
three subsets rather than a fourth collection.}
\label{tab:dataset-composition}
\begin{tabular}{@{}lrrrr@{}}
\toprule
& Graph ML & Computer vision & Time series & All \\
\midrule
Ideas & 5{,}225 & 8{,}291 & 7{,}595 & 21{,}111 \\
Pairs & 21{,}183 & 36{,}238 & 35{,}455 & 92{,}876 \\
\addlinespace
\textsc{low} & 819 & 901 & 1{,}724 & 3{,}444 \\
\textsc{weak} & 1{,}240 & 1{,}711 & 1{,}614 & 4{,}565 \\
\textsc{moderate} & 2{,}247 & 2{,}777 & 2{,}173 & 7{,}197 \\
\textsc{high} & 919 & 2{,}902 & 2{,}084 & 5{,}905 \\
\bottomrule
\end{tabular}
\end{table}

\paragraph{Splits.}
Splits are date-ordered rather than random, so that no training target is dated
after an evaluation target. The cuts fall at 2025-02 and 2025-11, giving
$13{,}815/3{,}944/3{,}352$ ideas for training, validation, and test.

\subsection{Validation of Research Novelty Labels}
\label{app:dataset-validity}

Our novelty labels are intended to describe the contribution remaining beyond
retrieved prior work. We examine this interpretation in two ways: supervised
probes predict the stored labels from alternative inputs, and an exploratory
blind re-assessment tests responses to controlled changes in evidence.

\paragraph{Prediction from metadata and coverage features.}
We first measure associations with the continuous novelty score, recomputed
from the stored annotation components using Sec.~\ref{sec:dataset-task}.
Numeric features are target character length, word count, publication date,
number of retained priors, mean target--prior semantic
similarity, and mean prior coverage score. Categorical features are innovation
aspect, contribution type, and task family. We use Spearman $\rho$ for numeric
features and Kruskal--Wallis $\varepsilon^2$ for categorical features, with
reporting thresholds $|\rho|\ge0.20$ and $\varepsilon^2\ge0.06$.
Text length, word count, and publication date have weak associations
($\rho=+0.049$, $+0.058$, and $-0.130$); all three categorical effects are
$\varepsilon^2\le0.033$. Prior count, mean semantic similarity, and mean
coverage have negative associations ($\rho=-0.286$, $-0.232$, and $-0.471$).

To test their combined predictive value, we fit two class-balanced logistic
regression classifiers to the stored four-class novelty labels. The first uses
all features listed above except mean coverage; the second adds mean coverage
to that same feature set.

Quadratic weighted kappa (QWK) against the stored labels rises from $0.227$ without coverage to $0.539$ with coverage. This
comparison shows the additional predictive value of the coverage annotation
beyond the other recorded features.

\paragraph{Prediction with retrieved and random priors.}
We next train a text probe to test whether the choice of prior content affects
prediction of the same stored labels. In the \textsc{true} condition, the input
concatenates the target text and up to five retained prior texts. In the
\textsc{random} condition, it concatenates the same target with up to five
priors sampled uniformly without replacement from the deduplicated prior pool,
restricted to dates earlier than the target.

The probe is TF--IDF followed by class-balanced logistic regression. TF--IDF
uses lowercased unigrams and bigrams, sublinear term frequency, a minimum
document frequency of two, and at most $200{,}000$ features.

Test QWK is $0.196$ with retrieved priors and $0.005\pm0.021$ with random
priors (mean and standard deviation across seeds), a gap of $0.192$.
Retrieved prior content therefore carries more information about the stored
labels than random earlier content under this probe.

\paragraph{Response to evidence covering the target contribution.}
The prediction experiments keep the labels fixed. We now test whether novelty
judgments respond appropriately when the evidence changes. If an earlier study
already contains the target contribution, the assigned novelty should decrease.
We sample $30$ targets across the four grades and construct one synthetic prior
for each target by copying its technical text and assigning a date one month
earlier. After adding it to the evidence set, we use \texttt{deepseek-flash}
to perform a blind re-assessment of pair coverage, joint coverage, residual
contribution, evidence sufficiency, and surface-level difference. The judge sees
only the texts, dates, and anonymized prior identifiers; the novelty score and
grade are computed using the rule of Sec.~\ref{sec:dataset-task}. Synthetic
priors have dedicated identifiers and are marked \texttt{synthetic\_control}
in the experiment records. They do not refer to real papers, authors, or DOIs
and are excluded from the released dataset.

Each modified target is compared with the medians of five independent blind
assessments of its original evidence set to account for variation across
repeated assessments. We evaluate recognition of the added prior as large
coverage, changes in joint coverage and residual contribution, and decreases
in the novelty score and grade. Ties are excluded from the exact two-sided
sign tests.

The added prior is labeled large coverage in all $29$ comparable cases.
Joint coverage increases and residual contribution decreases in $27$ cases
each (sign test $p=2.2\times10^{-7}$ for each component). The novelty score
decreases in $27/29$ cases ($p=1.6\times10^{-6}$), and the novelty grade
decreases for $82.8\%$ of targets. These results show that the blind judge
lowers novelty when the supplied evidence covers the target contribution.

\paragraph{Stability under changes that preserve the contribution.}
A novelty judgment should also remain stable when the evidence changes without
altering the substantive contribution of the target. For the same $30$ targets,
we reorder the priors, duplicate existing priors, or add unrelated papers.
The novelty grades before and after these changes show agreement measured by
quadratic weighted Cohen's $\kappa$ \citep{cohen1960kappa} of $0.710$, $0.812$,
and $0.889$, respectively. These results complement the coverage experiment:
the judgments respond to evidence that reduces the remaining contribution and
remain comparatively stable under changes that preserve it.
